\section{Methodology}\label{coll}

Our identification strategy leverages a Regression Discontinuity Design (RDD) to create a quasi-experimental setting that compares cases where the incumbent narrowly wins a contract with those where it barely loses. We focus on auctions where the bid difference between the winning and the next-best offer is minimal. In such tight contests, under competitive conditions, the outcome—whether a firm wins or loses—should be essentially random, as the firms involved are nearly identical in observable and unobservable characteristics. Any abrupt change in outcomes at this narrow margin may therefore be attributed to non-competitive factors, such as collusion or favoritism.

We define the running variable, $\Delta_{i,t}$, as the difference between firm $i$’s bid and the lowest competing bid in auction $t$. Negative values indicate that firm $i$ underbid its nearest competitor and won the auction, while positive values denote a loss. The threshold is set at $\Delta=0$, which distinguishes between narrowly winning and narrowly losing firms. This setting allows us to compare auctions that are nearly tied, under the assumption that there is no systematic difference between the two groups except for the outcome of winning.

The key to our approach is the assumption that firms cannot precisely manipulate their bids to land exactly near the cutoff. In a sealed-bid environment, where bids are submitted without knowledge of rivals’ offers, such precise manipulation is unlikely. Consequently, near-tie auctions provide a naturally randomized comparison: if the bidding were purely competitive, the probability of winning in these close cases should be approximately 50\%, and firm characteristics should vary smoothly across the cutoff.

We employ a local linear regression to estimate the causal effect of an incumbent win. For each auction $i$, we model the outcome $Y_{i}$ (for example, the firm’s backlog, incumbency status, or recent win rate) as a function of $\Delta_{i,t}$:
\[
Y_{i} \;=\; \alpha_0 + \tau \cdot \mathbf{1}\{\Delta_{i,t} < 0\} + \beta_1 f(\Delta_{i,t}) + \beta_2 f(\Delta_{i,t}) \cdot \mathbf{1}\{\Delta_{i,t} < 0\} + \varepsilon_{i},
\]
where $\mathbf{1}\{\Delta_{i,t} < 0\}$ is an indicator that equals one if firm $i$ wins (i.e., its bid is lower than the closest competitor’s bid), and $f(\Delta_{i,t})$ represents a smooth function of the running variable (typically a first-order polynomial). The parameter $\tau$ captures the discontinuity at the cutoff and thus the causal effect of winning on the outcome. We select the bandwidth optimally using the method of Calonico, Cattaneo, and Titiunik (2014), which balances bias and variance. Robustness is further ensured by conducting sensitivity analyses with different bandwidths and polynomial orders.

To validate our identification, we verify that covariates such as the number of bidders and reserve prices change smoothly around $\Delta=0$. We also apply a McCrary density test to ensure that there is no abnormal clustering of observations at the cutoff, which would suggest strategic manipulation. The absence of such manipulation reinforces the credibility of our RDD design.

Our analysis focuses on several outcome variables. Chief among these is the \emph{incumbency status} of the winning firm—whether it had won the previous similar auction—and its \emph{backlog}, defined as the standardized total value of ongoing contracts in the months prior to the auction. Additionally, we examine the \emph{recent win rate} of firms, which is the share of past auctions a firm has won within a specified period. In a competitive setting, these outcomes should not differ significantly between nearly tied auctions. However, if collusion (particularly bid rotation) is at work, we expect to see systematic differences: for instance, narrowly winning firms may exhibit lower recent win rates or smaller backlogs, consistent with a rotation scheme designed to distribute wins among cartel members.

We also extend our analysis to study the impact on winning bid prices. Collusion is often associated with higher contract prices due to the lack of genuine competition. Thus, by examining the ratio of the winning bid to the reserve price (or an analogous price index), we can assess whether auctions with collusive signals also exhibit elevated prices.

In summary, our RDD-based methodology provides a rigorous framework to isolate the causal effect of narrowly winning an auction on various outcomes. By focusing on close-margin contests, validating the continuity of covariates around the cutoff, and employing robust estimation techniques, we can credibly attribute any observed discontinuities to the treatment effect of incumbent win, shedding light on the underlying dynamics of collusion in public procurement. The following sections present our empirical findings, which build on this methodological foundation to reveal distinct patterns of anti-competitive behavior in the data.
